\documentclass[10pt, aspectratio=169]{beamer}
% Preámbulo modular para presentaciones Beamer (Modelación Estadística)
% Diseño y paleta institucional del Tecnológico de Monterrey

\documentclass[aspectratio=169,xcolor={dvipsnames,table}]{beamer}

\usepackage[utf8]{inputenc}
\usepackage[spanish,mexico]{babel}
\usepackage{amsmath,amssymb,amsfonts,mathtools}
\usepackage{graphicx}
\usepackage{listings}
\usepackage{booktabs}
\usepackage{tikz}
\usetikzlibrary{matrix,arrows,decorations.pathmorphing,shapes.geometric,positioning}

% Paleta Institucional Tecnológico de Monterrey
\definecolor{TecAzul}{HTML}{0039A6}       % Azul TEC: color primario institucional
\definecolor{TecAzulOscuro}{HTML}{1A2E51} % Azul marino oscuro
\definecolor{TecRojo}{HTML}{EC2661}       % Rojo de acento (paleta miTec)
\definecolor{TecGris}{HTML}{646464}       % Gris neutro para comentarios y subtítulos
\definecolor{TecFondoBloque}{HTML}{F4F6F9}% Fondo suave para bloques

% Configuración de Tema Beamer
\usetheme{Boadilla}
\usecolortheme[named=TecAzulOscuro]{structure}
\setbeamercolor{alerted text}{fg=TecRojo}
\setbeamercolor{palette primary}{bg=TecAzulOscuro,fg=white}
\setbeamercolor{palette secondary}{bg=TecAzul,fg=white}
\setbeamercolor{palette tertiary}{bg=TecAzulOscuro!80!black,fg=white}
\setbeamercolor{title}{fg=TecAzulOscuro,bg=TecFondoBloque}
\setbeamercolor{frametitle}{fg=TecAzulOscuro,bg=TecFondoBloque}

% Bloques personalizados elegantes
\setbeamercolor{block title}{bg=TecAzulOscuro,fg=white}
\setbeamercolor{block body}{bg=TecFondoBloque,fg=black}
\setbeamercolor{block title alerted}{bg=TecRojo,fg=white}
\setbeamercolor{block body alerted}{bg=TecRojo!8,fg=black}
\setbeamercolor{block title example}{bg=TecAzul,fg=white}
\setbeamercolor{block body example}{bg=TecAzul!8,fg=black}

% Redondear bordes y añadir sombra sutil
\setbeamertemplate{blocks}[rounded][shadow=true]
\setbeamertemplate{navigation symbols}{}

% Pie de página limpio con número de diapositiva
\setbeamertemplate{footline}{
  \leavevmode%
  \hbox{%
  \begin{beamercolorbox}[wd=.7\paperwidth,ht=2.25ex,dp=1ex,left]{author in head/foot}%
    \hspace*{2ex}\insertshortauthor~~(\insertshortinstitute) \hfill \insertshorttitle
  \end{beamercolorbox}%
  \begin{beamercolorbox}[wd=.3\paperwidth,ht=2.25ex,dp=1ex,right]{date in head/foot}%
    \insertshortdate{}\hspace*{2em}
    \insertframenumber{} / \inserttotalframenumber\hspace*{2ex}
  \end{beamercolorbox}}%
  \vskip0pt%
}

% Configuración de Listings de Python para visualización pedagógica de código
\lstset{
    language=Python,
    basicstyle=\ttfamily\footnotesize,
    keywordstyle=\color{TecAzulOscuro}\bfseries,
    stringstyle=\color{TecRojo},
    commentstyle=\color{TecGris}\itshape,
    showstringspaces=false,
    breaklines=true,
    frame=lines,
    rulecolor=\color{TecAzulOscuro!30},
    backgroundcolor=\color{TecFondoBloque},
    aboveskip=0.5em,
    belowskip=0.5em,
    literate={á}{{\'a}}1 {é}{{\'e}}1 {í}{{\'i}}1 {ó}{{\'o}}1 {ú}{{\'u}}1
             {Á}{{\'A}}1 {É}{{\'E}}1 {Í}{{\'I}}1 {Ó}{{\'O}}1 {Ú}{{\'U}}1
             {ñ}{{\~n}}1 {Ñ}{{\~N}}1 {¿}{{?`}}1 {¡}{{!`}}1
}

% Metadatos institucionales por defecto
\title[Modelación Estadística]{Modelación Estadística para Ciencia de Datos e Ingeniería}
\author[J. Castillo Colmenares]{Juliho David Castillo Colmenares}
\institute[Tec de Monterrey]{Tecnológico de Monterrey \\ Escuela de Ingeniería y Ciencias}
\date{\today}

% Comandos y macros estadísticas compartidas para las presentaciones Beamer (Metropolis)

% Conjuntos numéricos básicos
\newcommand{\R}{\mathbb{R}}
\newcommand{\N}{\mathbb{N}}
\newcommand{\Q}{\mathbb{Q}}
\newcommand{\Z}{\mathbb{Z}}
\newcommand{\set}[1]{\left\{ #1 \right\}}
\newcommand{\abs}[1]{\left|#1\right|}
\newcommand{\norm}[1]{\left\| #1 \right\|}

% Comandos estadísticos y probabilísticos del curso
\newcommand{\Var}{\operatorname{Var}}
\newcommand{\cov}{\operatorname{Cov}}
\newcommand{\corr}{\rho}
\newcommand{\comb}[2]{\begin{pmatrix} #1 \\ #2 \end{pmatrix}}
\newcommand{\s}{\sigma}
\newcommand{\particion}{\mathfrak{P}}
\newcommand{\card}[1]{\#\left( #1 \right)}
\newcommand{\seq}[3]{\set{#1}_{#2}^{#3}}
\newcommand{\E}{\mathbb{E}}
\newcommand{\Prob}{\mathbb{P}}

% Entornos pedagógicos y cajas en español académico natural
\newenvironment{discusion}{%
  \begin{alertblock}{Pregunta de Discusión}%
}{%
  \end{alertblock}%
}

\newenvironment{reflexion}{%
  \begin{alertblock}{Reflexión para el Aula}%
}{%
  \end{alertblock}%
}

\newenvironment{paradoja}{%
  \begin{alertblock}{Paradoja de Exploración}%
}{%
  \end{alertblock}%
}

\newenvironment{motivacion}{%
  \begin{exampleblock}{Motivación: Ciencia de Datos en la Práctica}%
}{%
  \end{exampleblock}%
}

\newenvironment{retoclase}{%
  \begin{block}{Reto Rápido en Equipo}%
}{%
  \end{block}%
}


\title[Esperanza y Varianza Continua]{Sección 04.03: Esperanza y Varianza Continua}
\subtitle{Definición Integral, LOTUS y Descomposición de Varianza}
\author[J. Castillo Colmenares]{Juliho Castillo Colmenares}
\institute{Tecnológico de Monterrey}
\date{\vspace{-1.2cm}}

\begin{document}

% Slide 1: Portada
\begin{frame}[plain]
  \titlepage
\end{frame}

% Slide 2: Hoja de Ruta
\begin{frame}{Hoja de Ruta --- Capítulo 04: Variables Aleatorias Continuas}
  \footnotesize
  ¿Dónde nos encontramos en el desarrollo modular del capítulo? \pause
  \begin{itemize}\itemsep=0.08em
    \item \textbf{04.01} Función de Densidad (PDF) y Soporte Continuo \pause
    \item \textbf{04.02} Función de Distribución Acumulada (CDF) Continua y Cuantiles \pause
    \item \textbf{04.03} Esperanza, Varianza y LOTUS Continuo \emph{(Hoy)} \pause
    \item \textbf{04.04} Distribución Uniforme Continua \pause
    \item \textbf{04.05} Distribución Normal y Puntaje $Z$ \pause
    \item \textbf{04.06} Distribución Exponencial y Procesos sin Memoria \pause
    \item \textbf{04.07} Distribuciones Gamma, Beta y Weibull
  \end{itemize}
  \vspace{-0.05cm}
  \begin{block}{Objetivo de la Sesión}
    Definir la esperanza continua $\E[X] = \int x f(x)\,dx$ y la varianza $\Var(X) = \int (x-\mu)^{2} f(x)\,dx$, demostrar el Teorema LOTUS para transformaciones, y aplicar la ley total de varianza al análisis de modelos jerárquicos.
  \end{block}
\end{frame}

% Slide 3: Motivación
\begin{frame}{De la PDF a los Momentos: Integración como Operación Natural}
  \footnotesize
  \begin{motivacion}[¿Por qué la esperanza continua?]
    Conocemos la PDF y la CDF de una variable continua. Ahora necesitamos resúmenes numéricos que capturen las propiedades esenciales: la \emph{esperanza} $\E[X]$ (centro de masa de la distribución) y la \emph{varianza} $\Var(X)$ (dispersión alrededor del centro). Estos momentos guían decisiones en inferencia, simulación y análisis de riesgos.
  \end{motivacion}

  \vspace{0.15cm}
  \pause
  \begin{itemize}\itemsep=0.1cm
    \item \textbf{Esperanza:} $\E[X] = \int x f(x)\,dx$, generalización natural de $\E[X] = \sum x P(X=x)$ mediante sumas de Riemann. \pause
    \item \textbf{Varianza:} $\Var(X) = \E[(X - \mu)^{2}] = \E[X^{2}] - (\E[X])^{2}$, mide la dispersión cuadrática media. \pause
    \item \textbf{LOTUS:} $\E[g(X)] = \int g(x) f(x)\,dx$ evita derivar la distribución de $g(X)$.
  \end{itemize}
\end{frame}

% Slide 4: Esperanza Continua
\begin{frame}{Esperanza Continua: $\E[X] = \int x f(x)\,dx$}
  \footnotesize
  Para una variable aleatoria continua $X$ con PDF $f(x)$, la \emph{esperanza matemática} (o media) se define como: \pause

  \vspace{0.1cm}
  \begin{block}{Definición de Esperanza Continua}
    $$\E[X] = \int_{-\infty}^{\infty} x \cdot f(x)\,dx$$
    bajo la condición de convergencia $\int |x| f(x)\,dx < \infty$. Para distribuciones de colas ligeras (Exponencial, Normal, Uniforme), la integral converge; para colas pesadas (Cauchy), diverge.
  \end{block}

  \vspace{0.1cm}
  \pause
  \begin{alertblock}{Interpretación Física}
    \begin{itemize}\itemsep=0.05cm
      \item $\E[X]$ es el \emph{centro de masa} de la distribución: equilibra la PDF como una balanza.
      \item Generaliza la media aritmética: si $X$ es uniforme en $[a, b]$, entonces $\E[X] = (a+b)/2$.
      \item La esperanza no necesita estar en el soporte: para PDF triangular en $[0, 1]$, $\E[X] = 2/3 \in [0, 1]$ \checkmark.
    \end{itemize}
  \end{alertblock}
\end{frame}

% Slide 5: Varianza Continua
\begin{frame}{Varianza Continua: $\Var(X) = \E[X^{2}] - (\E[X])^{2}$}
  \footnotesize
  La \emph{varianza} mide la dispersión cuadrática media de $X$ alrededor de su esperanza: \pause

  \vspace{0.1cm}
  \begin{block}{Definición de Varianza Continua}
    $$\Var(X) = \E[(X - \mu)^{2}] = \int_{-\infty}^{\infty} (x - \mu)^{2} f(x)\,dx = \E[X^{2}] - (\E[X])^{2}$$
    donde $\E[X^{2}] = \int x^{2} f(x)\,dx$ es el segundo momento bruto. La desviación estándar es $\sigma = \sqrt{\Var(X)}$.
  \end{block}

  \vspace{0.1cm}
  \pause
  \begin{alertblock}{Propiedades Fundamentales}
    \begin{itemize}\itemsep=0.05cm
      \item $\Var(X) \ge 0$ siempre, con igualdad si y solo si $X$ es constante.
      \item $\Var(aX + b) = a^{2} \Var(X)$ (linealidad afín).
      \item $\Var(X) = 0$ implica que $X$ es determinista; distribuciones concentradas tienen varianza pequeña.
    \end{itemize}
  \end{alertblock}
\end{frame}

% Slide 6: Teorema LOTUS
\begin{frame}{Teorema LOTUS: $\E[g(X)] = \int g(x) f(x)\,dx$}
  \footnotesize
  El Teorema LOTUS (Law of the Unconscious Statistician) permite calcular la esperanza de $g(X)$ sin necesidad de encontrar la distribución de $g(X)$: \pause

  \vspace{0.1cm}
  \begin{block}{LOTUS para Variables Continuas}
    Para cualquier función medible $g: \mathbb{R} \to \mathbb{R}$ y variable continua $X$ con PDF $f(x)$:
    $$\E[g(X)] = \int_{-\infty}^{\infty} g(x) \cdot f(x)\,dx$$
    bajo condiciones de regularidad ($g$ medible, $g(X)$ integrable). Caso particular: $g(x) = x$ recupera $\E[X]$; $g(x) = x^{2}$ recupera $\E[X^{2}]$.
  \end{block}

  \vspace{0.1cm}
  \pause
  \begin{alertblock}{Aplicaciones Prácticas}
    \begin{itemize}\itemsep=0.05cm
      \item Transformaciones: $\E[\log X]$, $\E[X^{2}]$, $\E[\sqrt{X}]$, $\E[\max(X - K, 0)^{+}]$ para opciones.
      \item Momentos superiores: asimetría $\gamma_{1} = \E[(X-\mu)^{3}]/\sigma^{3}$ y curtosis $\gamma_{2} = \E[(X-\mu)^{4}]/\sigma^{4} - 3$.
      \item Esperanza de indicadores: $\E[\mathbb{1}_{X > c}] = P(X > c) = 1 - F(c)$.
    \end{itemize}
  \end{alertblock}
\end{frame}

% Slide 7: Linealidad de la Esperanza
\begin{frame}{Linealidad de la Esperanza: $\E[aX + b] = a\E[X] + b$}
  \footnotesize
  La esperanza es un operador lineal, válido incluso para variables aleatorias \emph{dependientes}: \pause

  \vspace{0.1cm}
  \begin{block}{Teorema de Linealidad}
    \begin{itemize}\itemsep=0.05cm
      \item Para constantes $a, b$ y variable aleatoria $X$: $\E[aX + b] = a \E[X] + b$.
      \item Para variables aleatorias $X, Y$ (independientes o no): $\E[X + Y] = \E[X] + \E[Y]$.
      \item La demostración usa la propiedad de linealidad de la integral.
    \end{itemize}
  \end{block}

  \vspace{0.1cm}
  \pause
  \begin{alertblock}{Contraste con la Varianza}
    \begin{itemize}\itemsep=0.05cm
      \item $\Var(X + Y) = \Var(X) + \Var(Y) + 2\cov(X, Y)$, \emph{no} $\Var(X) + \Var(Y)$ en general.
      \item $\Var(X + Y) = \Var(X) + \Var(Y)$ solo si $X, Y$ son \emph{independientes} o \emph{incorrelacionadas}.
      \item La esperanza es robusta a la dependencia; la varianza requiere más estructura.
    \end{itemize}
  \end{alertblock}
\end{frame}

% Slide 8: Ley Total de Varianza
\begin{frame}{Ley Total de Varianza: Descomposición ANOVA}
  \footnotesize
  Para variables jerárquicas, la varianza total se descompone en variabilidad dentro y entre grupos: \pause

  \vspace{0.1cm}
  \begin{block}{Ley Total de Varianza}
    Si $X$ es una variable con $Y$ como índice de grupo, entonces:
    $$\Var(X) = \E[\Var(X \mid Y)] + \Var(\E[X \mid Y])$$
    El primer término es la \emph{variabilidad dentro de grupos}; el segundo es la \emph{variabilidad entre grupos}.
  \end{block}

  \vspace{0.1cm}
  \pause
  \begin{alertblock}{Aplicación en ANOVA}
    \begin{itemize}\itemsep=0.05cm
      \item El estadístico $F$ compara la variabilidad entre grupos con la variabilidad dentro de grupos.
      \item Si $F$ es grande, la mayor parte de la varianza es \emph{entre} grupos (tratamiento efectivo).
      \item La identidad se demuestra expandiendo $\E[(X - \E[X])^{2}]$ vía la ley total de esperanza.
    \end{itemize}
  \end{alertblock}
\end{frame}

% Slide 9: Aplicaciones
\begin{frame}{Aplicaciones en Ingeniería, Ciencias y Ciencia de Datos}
  \footnotesize
  Los momentos continuos (esperanza, varianza, asimetría) son herramientas centrales en análisis cuantitativo: \pause

  \vspace{0.15cm}
  \begin{itemize}\itemsep=0.12cm
    \item \textbf{Propagación de Incertidumbre:} Si $X = \mu + \epsilon$ con $\epsilon \sim N(0, \sigma^{2})$, el error estándar del promedio de $n$ mediciones es $\sigma/\sqrt{n}$. \pause
    \item \textbf{Cuantificación de Riesgo:} VaR al $95\%$ es el cuantil $q_{0.05}$ de la distribución de pérdidas; CVaR es la esperanza condicional $\E[X \mid X \le q_{0.05}]$. \pause
    \item \textbf{Análisis de Regresión:} En regresión lineal, los residuos deben tener esperanza cero y varianza constante (homocedasticidad). \pause
    \item \textbf{Control de Calidad:} Cartas de control $\bar{x}$ usan $\mu \pm 3\sigma/\sqrt{n}$ como límites de control.
  \end{itemize}
\end{frame}

% Slide 12A-1: Lab Python (Bloque 1)
\begin{frame}[fragile]{Laboratorio en Python: Esperanza y LOTUS}
  \scriptsize
  Cálculo de $\E[X]$, $\E[X^2]$ y verificación de LOTUS para $g(X) = \sqrt{X}$ y $g(X) = \log(X)$ (\texttt{04.03\_expectation\_and\_variance.py}):
  \vspace{0.02cm}
  {\lstset{basicstyle=\fontsize{4.5pt}{5.5pt}\selectfont\ttfamily}
  \lstinputlisting[firstline=15, lastline=38]{../../code/04_variables_aleatorias_continuas/04.03_expectation_and_variance.py}}
\end{frame}

% Slide 12A-2: Lab Python (Bloque 2)
\begin{frame}[fragile]{Laboratorio en Python: Varianza y Momentos Centrales}
  \scriptsize
  Cálculo de varianza y momentos centrales (asimetría, curtosis) para Uniform, Exponential, Normal, y demostración de la ley total de varianza:
  \vspace{0.02cm}
  {\lstset{basicstyle=\fontsize{4pt}{5pt}\selectfont\ttfamily}
  \lstinputlisting[firstline=45, lastline=75]{../../code/04_variables_aleatorias_continuas/04.03_expectation_and_variance.py}}
\end{frame}

% Slide 12B: Lab Python (Bloque 3)
\begin{frame}[fragile]{Laboratorio en Python: Monte Carlo y Propagación}
  \scriptsize
  Simulación Monte Carlo con $N = 250{,}000$ y análisis de propagación de error en medición:
  \vspace{0.02cm}
  {\lstset{basicstyle=\fontsize{4.5pt}{5.5pt}\selectfont\ttfamily}
  \lstinputlisting[firstline=85, lastline=108]{../../code/04_variables_aleatorias_continuas/04.03_expectation_and_variance.py}}
\end{frame}

% Slide 12C: Salida Terminal
\begin{frame}[fragile]{Laboratorio en Python: Salida en Terminal}
  \scriptsize
  Salida estándar generada al ejecutar el script del laboratorio en Python:
  \vspace{0.02cm}
  \begin{block}{Salida Estándar en Consola}
    \fontsize{4pt}{5pt}\selectfont
    \begin{verbatim}
=== Block 1: Expectation & LOTUS ===
Triangular(0,1): E[X]=0.666667 | E[X^2]=0.500000 | Var=0.055556
LOTUS E[sqrt(X)]=0.800000 | LOTUS E[log(X)]=-0.499993
Linearity E[3.5X+10]=12.333333

=== Block 2: Variance & Central Moments ===
Uniform(0,1): E[X]=0.500000 | Var=0.083333
Exponential(1.0): Skewness=2.000000 | Excess kurtosis=6.000000
Hierarchical: total Var=4.2187 = within 1.0221 + between 3.2991

=== Block 3: Monte Carlo & Error Propagation ===
Exponential MC (N=250,000): mean=1.0011 | var=1.0003
  Skewness=1.9909 | Kurtosis=5.8731
Error propagation (sigma=0.1, true=5.0):
  n=1024: SE=0.0031 | CI half-width=0.0061
  Required n for SE<0.01: 100
    \end{verbatim}
  \end{block}
\end{frame}

% Slide 13A: Ejercicio Nivel 1 (Enunciado)

% Slide 13B: Ejercicio Nivel 1 (Resolución)

% Slide 14A: Ejercicio Nivel 2 (Enunciado)

% Slide 14B: Ejercicio Nivel 2 (Resolución)

% Slide 15A: Ejercicio Nivel 3 (Enunciado)

% Slide 17: Cierre
\begin{frame}{Cierre de Sección y Síntesis sobre Esperanza Continua}
  \begin{reflexion}[Los Momentos como Resumen Distributional]
    La esperanza continua $\E[X] = \int x f(x)\,dx$ y la varianza $\Var(X) = \int (x - \mu)^{2} f(x)\,dx$ resumen las propiedades esenciales de cualquier distribución. El Teorema LOTUS extiende estas ideas a transformaciones sin necesidad de derivar nuevas PDFs, y la ley total de varianza conecta momentos con estructuras jerárquicas.
  \end{reflexion}

  \vspace{0.2cm}
  \pause
  \begin{block}{Lecciones Clave de la Esperanza Continua}
    \begin{itemize}\itemsep=0.08cm
      \item \textbf{Esperanza:} $\E[X]$ es el centro de masa; la integral $\int x f(x)\,dx$ generaliza la suma discreta.
      \item \textbf{Varianza:} $\Var(X) = \E[X^{2}] - (\E[X])^{2}$; las varianzas suman solo si las variables son independientes.
      \item \textbf{LOTUS:} $\E[g(X)] = \int g(x) f(x)\,dx$ evita derivar la distribución de $g(X)$.
    \end{itemize}
  \end{block}
\end{frame}

% Slide 18: Perspectiva Modular
\begin{frame}{Perspectiva Modular --- El Próximo Reto}
  \scriptsize
  \begin{retoclase}[Para la próxima sesión: Distribución Uniforme Continua]
    Con la PDF, la CDF y los momentos en mano, podemos ahora estudiar distribuciones específicas. La más simple es la \emph{Uniforme Continua} $U(a, b)$, donde $f(x) = 1/(b-a)$ en $[a, b]$. Aunque elemental, es fundamental para generación de variables aleatorias, simulación Monte Carlo y el método de inversión.
  \end{retoclase}

  \vspace{0.08cm}
  \pause
  \begin{alertblock}{Preparación Recomendada}
    \begin{itemize}\itemsep=0.06cm
      \item Repasa el concepto de variable equiprobable y su generalización al continuo.
      \item Reflexiona sobre cómo la uniforme es la base para simulación de cualquier otra distribución vía transformación.
    \end{itemize}
  \end{alertblock}
\end{frame}

\end{document}
